% !TEX root = ../main.tex

\chapter{Entropia układów kwantowych}\label{sec:entropia_kwantowa}
W stanie równowagi termodynamicznego prawdopodobieństwa znalezienia układu w stanie o pewnej energii opisuje wzór Boltzmanna, co można zapisać w formie operatorowej jako \cite{b:huang_podstawy_fizyki_statystycznej}:
\begin{gather}
	\vu{\rho}^{eq} = \inv{Z} \exp\pab{-\beta\ham}, \\
	\ln\vu{\rho}^{eq} \overset{def}{=} - \ln Z - \beta\ham,
\end{gather}
gdzie \(Z = \Tr\Bab{\exp\pab{-\beta\ham}}\) to suma statystyczna. Odpowiednio jak w klasycznej fizyce statystycznej postulujemy, że energia swobodna Helmholtza \(F\) spełnia \cite{b:huang_mechanika_statystyczna}:
\begin{equation*}
	F = - \inv{\beta} \ln Z.
\end{equation*}
Prawdziwość tej równości możemy zweryfikować posługując się znanymi tożsamościami \cite{b:pudlik_termodynamika}:
\begin{gather*}
	F = U - TS = U - \frac{S}{k_B \beta}, \\
	\pab{\pdv{U}{S}}_V = T = \inv{k_B \beta}, \\
	\pab{\pdv{F}{T}}_V = -S, \\
	\pab{\pdv{F}{\beta}}_V = \pab{\pdv{F}{T}}_V \pdv{T}{\beta} = \frac{S}{k_B \beta^2}.
\end{gather*}
Wykorzystując je możemy wyprowadzić \eqref{eq:kwant_energia_wewnetrzna}:
\begin{multline*}
	U = F + \frac{S}{k_B \beta} =
	F + \beta \pab{\pdv{F}{\beta}}_V =
	\bab{\pdv{}{\beta}\pab{\beta F}}_V =
	- \bab{\pdv{}{\beta}\ln Z}_V = \\ =
	- \inv{Z} \pab{\pdv{Z}{\beta}}_V =
	- \inv{Z} \Tr\Bab{- \ham \exp\pab{-\beta\ham}} =
	\inv{Z} \Tr\Bab{Z \vu{\rho}^{eq} \ham} =
	\Tr\Bab{\vu{\rho}^{eq} \ham}.
\end{multline*}
Entropię można obliczyć przekształceniami:
\begin{multline*}
	S =
	k_B \beta \pab{U - F} =
	- k_B \beta F - k_B \beta U =
	k_B \ln Z - k_B \beta \Tr\Bab{\vu{\rho}^{eq} \ham} = \\ =
	k_B \frac{\ln Z}{Z} \Tr\Bab{\exp\pab{-\beta\ham}} + k_B \beta\Tr\Bab{\inv{Z} \exp\pab{-\beta\ham} \ham} = \\ =
	- k_B \Tr\Bab{\inv{Z}\exp\pab{-\beta\ham} \bab{- \ln Z - \beta\ham}} =
	- k_B \Tr\Bab{\vu{\rho}^{eq} \ln \vu{\rho}^{eq}}.
\end{multline*}